\subsubsection*{命题 6.1.}
$\tau^{-1}$ 线性运算 $\bar\alpha_S:W_S(1)\to W_S(1)$ 等于

\begin{equation}\tag{6.1.1}
\frac1p\,\tau^{-1}\circ\psi_p\circ(X_0f(X))^{p-1}.
\end{equation}

\noindent\emph{证明.}
写作

\begin{equation}\tag{6.1.2}
H(X)=1+\sum_{i\geq1}H_i(X),
\end{equation}

其中

\begin{equation}\tag{6.1.3}
H_i(X)\text{ 关于下述变量的次数为 $i$: }X_0.
\end{equation}

再次由 (4.6.19),

\begin{equation}\tag{6.1.4}
\frac1p\,\psi_pR_S(i)\subset pR_S
\qquad\text{若 }i\ne p.
\end{equation}

由于

\begin{equation}\tag{6.1.5}
H_i(X)\,R_S(j)\subset R_S(i+j),
\end{equation}

有

\begin{equation}\tag{6.1.6}
\left(
\alpha_S-\frac1p\,\tau^{-1}\circ\psi_p\circ H_{p-1}(X)
\right)(R_S(1))
\subset pR_S.
\end{equation}

因此,

\begin{equation}\tag{6.1.7}
\left[
(\alpha_S)^{\operatorname{ord}(q)}
-
\left(
\frac1p\,\tau^{-1}\circ\psi_p\circ H_{p-1}(X)
\right)^{\operatorname{ord}(q)}
\right]R_S(1)
\subset pR_S.
\end{equation}

尚需计算 $H_{p-1}(X)$. 由 (4.4.4),

\begin{equation}\tag{6.1.8}
\begin{aligned}
H(X)
&=\prod_W\Theta(\pi A_WX_0X^W)\\
&=\prod_W\exp\!\left(
\sum_{r\geq0}\frac{(\pi A_WX_0X^W)^{p^r}}{p^r}
\right)\\
&=\prod_W\prod_{r\geq0}
\exp\!\left(\frac{(\pi A_WX_0X^W)^{p^r}}{p^r}\right)\\
&\equiv\prod_W\exp(\pi A_WX_0X^W)
\pmod{\text{次数为如下范围的项: }\geq p}\\
&=\exp\!\left(\pi\sum_WA_WX_0X^W\right)
\pmod{\text{次数为如下范围的项: }\geq p}\\
&=\exp(\pi X_0f(X))
\pmod{\text{次数为如下范围的项: }\geq p}.
\end{aligned}
\end{equation}

由此得

\begin{equation}\tag{6.1.9}
H_{p-1}(X)
=\frac{\pi^{p-1}X_0^{p-1}}{(p-1)!}\,f(X)^{p-1}.
\end{equation}

容易验证

\begin{equation}\tag{6.1.10}
\frac{\pi^{p-1}}p\equiv-1\pmod p,
\end{equation}

从而

\begin{equation}\tag{6.1.11}
\frac1p\,H_{p-1}(X)
\equiv(X_0f(X))^{p-1}\pmod p.
\end{equation}

这就完成了命题~6.1 的证明.

\subsubsection*{推论 6.1.12.}
线性算子
\[
(\bar\alpha_S)^{\operatorname{ord}(q)}:W_S(1)\longrightarrow W_S(1)
\]
恰为
\[
\psi_q\circ(X_0f(X))^{q-1}.
\]

\noindent\emph{证明.}
由 (6.1.1),

\begin{equation*}
\begin{aligned}
(\bar\alpha_S)^{\operatorname{ord}(q)}
&=
\left(
\frac1p\,\tau^{-1}\circ\psi_p\circ(X_0f(X))^{p-1}
\right)^{\operatorname{ord}(q)}\\
&=
\psi_q\circ
\prod_{j=0}^{\operatorname{ord}(q)-1}
\left(X_0^{p^j}f^{\tau^j}(X^{p^j})\right)^{p-1}\\
&=
\psi_q\circ
\prod_{j=0}^{\operatorname{ord}(q)-1}
(X_0f(X))^{p^j(p-1)}\\
&=\psi_q\circ(X_0f(X))^{q-1}.
\end{aligned}
\end{equation*}

\medskip
\noindent\emph{底本注.}
印本证明中的第一个乘积具有未定义的上界 $\operatorname{ord}(a)-1$;
下一行及 $q$ 的定义要求为 $\operatorname{ord}(q)-1$.
校订法文层和英文层采用后者; 照录法文层保留印本形式.

\subsubsection*{推论 6.1.13.}
模 $p$ 有同余式

\begin{equation*}
Z(t,V/\mathbb F_q)
\equiv
\frac{
\det\!\left(
1-t\,\psi_q\circ(X_0f(X))^{q-1}
\mathrel{\Big|}W_S(1)
\right)^{(-1)^{n+1}}
}{1-t}.
\end{equation*}

